% ======================================================================
% Appendix A -- Setting Up ADK and Google Cloud
% ======================================================================

\chapter{Setting Up ADK and Google Cloud}
\label{app:setup}

This appendix covers the one-time setup every lab assumes, and the cost controls and cleanup habits that keep a semester of labs affordable. Commands change; the companion repository's \texttt{setup/} folder holds the current versions of every script referred to here, and its README records the versions they were last tested with.

\section{Local tools}
\label{sec:app-tools}

\begin{enumerate}[itemsep=2pt]
  \item \textbf{Python.} Install Python 3.11 or later. The labs use \texttt{uv} to manage virtual environments and pinned dependencies; \texttt{pip} also works.
  \item \textbf{The companion repository.} Clone it and install the dependencies of the chapter you are working on:
\begin{lstlisting}[language=bash, numbers=none]
git clone <companion-repository-url> agentic-solutions
cd agentic-solutions/ch01
uv sync            # or: python -m venv .venv && pip install -r requirements.txt
\end{lstlisting}
  Each chapter pins its own versions of ADK\index{Agent Development Kit (ADK)}, the MCP\index{Model Context Protocol (MCP)} SDK and other libraries, so labs keep working when new versions are released.
  \item \textbf{Google Cloud CLI.} Install \texttt{gcloud} and sign in twice---once for the CLI and once for application default credentials, which ADK and the client libraries use:
\begin{lstlisting}[language=bash, numbers=none]
gcloud auth login
gcloud auth application-default login
\end{lstlisting}
  \item \textbf{Terraform} (optional), for labs that provision infrastructure from the samples' scripts.
\end{enumerate}

\section{A Google Cloud account and project for the course}
\label{sec:app-account}

\subsection{Recommended: each student uses the Google Cloud free trial}

The recommended way to run the labs is for each student to sign up for Google Cloud's free trial. New customers receive a \$300 credit to spend over 90 days, and the trial cannot produce a bill: if the credit is used up or the 90 days pass without the account being upgraded, the trial billing account closes and its projects' resources stop. For a course, this means every student has a private, capped budget, and the instructor does not need a billing account at all. Every lab in this book runs within the trial, with the adjustments described below, and a student who runs each lab's cleanup script should use only a fraction of the credit over the course. Check the current terms on Google's free-program pages before the course starts; the details below were correct when this book went to press.

\paragraph{Signing up.}
\begin{enumerate}[itemsep=2pt]
  \item Go to Google Cloud's free-trial page and choose to start free, using a personal Google account rather than one managed by the university, unless the university's Google administrator has enabled Google Cloud for student accounts.
  \item Provide a credit card or other payment method. Google uses it to verify identity; it is not charged during the trial. Students who cannot or prefer not to provide one should follow the alternatives at the end of this section.
  \item \textbf{Do not upgrade to a paid account.} Upgrading removes the automatic stop at \$300 and 90 days, and spending beyond the credit is then billed.
\end{enumerate}

\paragraph{When to sign up.} The 90 days start at signup and cannot be paused or extended. A 14-week semester is 98 days, so a student who signs up in week~1 loses the credit before the capstone presentations. Sign up \textbf{at the start of week~2 or~3}, and complete the local-only parts of Labs~1--3 before then. In a 10-week course, sign up in week~1.

\paragraph{Creating the course project.} With the trial active, create one project for the course and link it to the trial billing account. You need the account's ID, which is different for every student. List your billing accounts:

\begin{lstlisting}[language=bash, numbers=none]
gcloud billing accounts list
\end{lstlisting}

The output looks like this:

\begin{lstlisting}[numbers=none]
ACCOUNT_ID            NAME                OPEN  MASTER_ACCOUNT_ID
01A2B3-C4D5E6-F7G8H9  My Billing Account  True
\end{lstlisting}

The Free Trial account is usually named \emph{My Billing Account}, and \texttt{OPEN} must be \texttt{True}. Copy its \texttt{ACCOUNT\_ID} into the variable \texttt{BILLING} below, replacing \texttt{XXXXXX-XXXXXX-XXXXXX}, then create the project, link it and enable the services the labs use:

\begin{lstlisting}[language=bash, numbers=none]
export BILLING=XXXXXX-XXXXXX-XXXXXX      # your ACCOUNT_ID
export PROJECT=agentic-course-$RANDOM
gcloud projects create $PROJECT
gcloud config set project $PROJECT
gcloud billing projects link $PROJECT --billing-account=$BILLING
gcloud services enable aiplatform.googleapis.com run.googleapis.com \
    pubsub.googleapis.com secretmanager.googleapis.com \
    cloudbuild.googleapis.com artifactregistry.googleapis.com \
    logging.googleapis.com monitoring.googleapis.com cloudtrace.googleapis.com
\end{lstlisting}

If you have only one open billing account, the shell can fill in the ID for you instead of the first line above:

\begin{lstlisting}[language=bash, numbers=none]
export BILLING=$(gcloud billing accounts list --filter=open=true \
    --format="value(name.basename())" --limit=1)
echo $BILLING                            # check it before linking
\end{lstlisting}

Do not paste \texttt{XXXXXX-XXXXXX-XXXXXX} as written; it is a placeholder, and the link command fails with it. These commands assume you have already signed in with both \texttt{gcloud auth login} and \texttt{gcloud auth application-default login} (Section~\ref{sec:app-tools}). Keep \texttt{BILLING} set in the same terminal: the budget command in the cost-controls section below uses it too.

The first service in the list provides access to Gemini models and Agent Platform\index{Gemini Enterprise Agent Platform} services (formerly Vertex AI). Service names occasionally change after product renames; if a command reports an unknown service, check the repository's \texttt{setup/apis.txt} for the current list. Labs that use search, sessions or other managed services enable those APIs in their own setup scripts.

Each lab's \texttt{.env.example} lists the variables it needs. Copy it to \texttt{.env} and set at least the project, the region and the model name. Never commit \texttt{.env}.

\subsection{What the trial covers, and the adjustments the labs need}

The trial credit pays for the services the labs use: Gemini models called through your Google Cloud project, Cloud Run, Agent Runtime, Agent Search, Pub/Sub, Eventarc, Cloud Scheduler, Secret Manager, Cloud Storage, Firestore and the observability services. Lab-sized use of Cloud Run, Pub/Sub, Cloud Storage, Cloud Build, Secret Manager and logging usually falls within Google Cloud's always-free monthly limits and does not touch the credit at all. Four restrictions of the trial matter for this course:

\begin{description}[leftmargin=1em,itemsep=3pt]
  \item[Gemini through your project, not through AI Studio.] The trial credit does not pay for the Gemini API when it is used with an AI Studio API key. The labs call Gemini through the course project (\texttt{GOOGLE\_GENAI\_USE\_VERTEXAI=True} with your project and region in \texttt{.env}), which is how the credit is used. After Lab~1, open the project's billing report and confirm that model usage appears against the credit.
  \item[Google's models only.] The credit does not pay for partner models offered as managed APIs. Keep every lab on Gemini; the other frameworks in Appendix~B are for reading and for later projects.
  \item[No GPUs.] Trial accounts cannot add GPUs to virtual machines. Only Lab~11 is affected: run its time-series foundation model on CPU, which is practical for the lab's 2{,}000 series. The chapter's nightly accelerator batch is a design discussion, not a lab step.
  \item[No quota increases.] Trial projects keep their default quotas, including model request rates. These are sufficient for every lab as written. For the larger evaluation runs---Chapter~\ref{ch:case-support}'s 600 replayed contacts, Chapter~\ref{ch:case-airline}'s storm replay---use the repository scripts' \texttt{--rate} option to throttle requests, or the smaller samples the labs specify.
\end{description}

Table~\ref{tab:trial-labs} summarizes the labs.

\begingroup
\small
\begin{xltabular}{\textwidth}{@{}l>{\raggedright\arraybackslash}X>{\raggedright\arraybackslash}p{4.4cm}@{}}
\caption{The labs on the free trial.}\label{tab:trial-labs}\\
\toprule
Lab & Main services & On the free trial \\
\midrule
\endfirsthead
\toprule
Lab & Main services & On the free trial \\
\midrule
\endhead
1--3 & Local ADK; Gemini; Python & Yes; Lab 3 needs no cloud \\
4 & Agent Search, Cloud Run, Eventarc & Yes \\
5 & Cloud Run, Agent Runtime & Yes; throttle the 150 adversarial scripts \\
6 & Pub/Sub, Cloud Run, IAP, Cloud Scheduler & Yes \\
7 & Cloud Run, Cloud Trace, Cloud Monitoring & Yes \\
8 & CPU model training, Cloud Run & Yes \\
9 & Agent Runtime, Firestore & Yes, with the 500-passenger replay \\
10 & Cloud Run, Secret Manager, Cloud Storage & Yes \\
11 & Time-series model batch & Yes, on CPU \\
12 & Agent Runtime, Cloud Run & Yes \\
Capstone & Varies & Yes, on Gemini and without GPUs \\
\bottomrule
\end{xltabular}
\endgroup

\subsection{Keeping the credit}

What uses up the credit is almost never a lab while it runs; it is a resource left running afterwards. Agent Runtime deployments, search data stores and databases cost money while they exist, whether or not anyone uses them. Three habits keep a semester well within \$300:
\begin{itemize}[itemsep=1pt]
  \item run each lab's \texttt{cleanup.sh} at the end of the lab, and check the console's resource list;
  \item look at the billing page once a week---instructors can make this a two-minute check at the start of a class;
  \item never upgrade the trial account.
\end{itemize}
If a trial ends early because the credit runs out, its projects stop, and the account enters a short grace period in which resources can be recovered by upgrading. Students should not upgrade simply to recover a lab; they should rebuild it from the repository, which is what the setup and cleanup scripts are for.

\subsection{If a student cannot use the free trial}

The trial is for new Google Cloud customers, so a student who has used it before cannot start another, and some students will not want to provide a payment method. Two alternatives keep them in the course:
\begin{itemize}[itemsep=2pt]
  \item \textbf{Local and free-tier route.} Labs~1--3 and the tool-design, evaluation and business parts of every other lab run on a laptop. For model calls, use a Gemini API key from Google AI Studio on its free tier, which needs no billing. Its rate limits are low, and Google uses content sent on the free tier to improve its products, which is acceptable here only because every lab uses synthetic data. Students deploy nothing, and submit the local version of each lab with its deployment section answered as a written plan.
  \item \textbf{A paid account with a budget.} A student with an existing billing account can use it, with a budget alert (next section), maximum-instance limits and the cleanup scripts. Most lab services stay within the always-free monthly limits; the paid exposure comes from Agent Runtime, Agent Search and model calls.
\end{itemize}
Instructors who prefer to fund the course centrally can apply to Google's credits-for-teaching program, through which faculty receive credits for their students; eligibility, amounts and countries are set by Google and change from time to time.

\section{Cost controls}

Set these before the first lab, not after the first surprise.
\begin{itemize}[itemsep=2pt]
  \item \textbf{A budget with alerts.} Create a budget with e-mail alerts at 50\%, 90\% and 100\% of a figure you choose---\$100 is a sensible ceiling for the whole course:
\begin{lstlisting}[language=bash, numbers=none]
gcloud billing budgets create --billing-account=$BILLING \
    --display-name="course budget" --budget-amount=100USD \
    --threshold-rule=percent=0.5 --threshold-rule=percent=0.9 \
    --threshold-rule=percent=1.0
\end{lstlisting}
  Here \texttt{\$BILLING} is the billing account ID from Section~\ref{sec:app-account}; set it again if you have opened a new terminal.
  A budget alerts; it does not stop spending. On the free trial, the trial's own \$300 cap is the hard stop; on a paid account, the next items are what limit spending.
  \item \textbf{Maximum instances} on every Cloud Run service (\texttt{--max-instances=2} is plenty for labs) and \textbf{minimum instances of zero}, so idle services cost nothing.
  \item \textbf{Agent budgets in code}: tool-call and model-call limits per request, as every case in this book uses. They are cost controls as much as safety controls.
  \item \textbf{Small models for development.} Use the smallest model that the lab's evaluation passes with; switch to larger ones only to compare.
  \item \textbf{No always-on resources.} Search indexes, databases and accelerators cost money while they exist, whether or not they are used. Create them when a lab needs them; delete them when it ends.
\end{itemize}

\section{Cleanup}

Every lab folder contains a \texttt{cleanup.sh} that deletes what the lab created: Cloud Run services and jobs, Pub/Sub topics and subscriptions, Agent Runtime deployments, search data stores, buckets and scheduled jobs. Run it at the end of each lab, then check the console's resource list for anything left behind. At the end of the course, deleting the whole project is the surest cleanup:
\begin{lstlisting}[language=bash, numbers=none]
gcloud projects delete $PROJECT
\end{lstlisting}

\section{When things go wrong}

\begin{description}[leftmargin=1em,itemsep=2pt]
  \item[Permission denied.] Check which identity is in use (\texttt{gcloud auth list}; application default credentials may differ from the CLI's) and whether the API is enabled. Deployed services run as their service account, not as you.
  \item[Model not found.] Model names change and older models are retired. Check the model name in \texttt{.env} against the repository's pinned value and the current model list for your region.
  \item[Works locally, fails deployed.] Usually a missing environment variable, a secret not granted to the service account, or a region mismatch between the service and the resources it calls.
  \item[An agent loops or runs up cost.] Look at the trace in \texttt{adk web} or Cloud Trace; check the tool-call limit; check whether a tool returns errors the model keeps retrying.
\end{description}
